Now we finally attempt to find the quantum theory. We only focus on
the case of the mainfold have topology $ M \cong \mathbb{R} \times
\Sigma_g $ for some genus $ g $ surface $ \Sigma_g $. We take the
spactime coordinate $ x^0 $ to be the $ \mathbb{R} $ direction and $
x^1, x^2 $ to be the spacial coordinates on $ \Sigma_g $.

As we'll see the
Lagrangian of the theory is a constraint system. Thus, we need to use
some quantization procedure for constraint systems. Apart from the
standard Dirac's procedure, here is more convenient to use another
procedure, called the "constraints first quantization".

\subsection{Constraint First Quantization}

\subsubsection{Classical Pre Constriant Phase Space}

Having a Lagrangian we can try to quantize the theory. We follow the
following steps:
\begin{enumerate}
  \item \textbf{First order in $ \partial_{0} $: } make sure all time
    derivative is in first order $
    \partial_{0} \phi $. If we have second order time derivative, we
    can introduce some auxiliary fields to make the Lagrangian first
    order. note that the Chern-Simons Lagrangian is already first order.
  \item \textbf{Separate out Dynamical Variables: } separate variable
    with a time derivative and those without.
    \begin{itemize}
      \item The ones with time derivative are the truly dynamical variables
      \item The ones without time derivative are the Lagrange
        multipliers, whose EoM gives us some constraints.
    \end{itemize}
  \item \textbf{Find pre-constraint phase space and Poisson Bracket:
    } find out the canonical variables and the phase space. Define a
    Poisson Bracket on the phase space.
\end{enumerate}

We then do it on the Lagrangian of 3D gravity with cosmological constant:
\begin{align}
  S = {} & -2 \int \mathrm{d}t \int_{\Sigma} \epsilon^{ij} \,
  e_{i\alpha} \partial_{0} \omega_{j}^{\alpha} \\
  & + \int \mathrm{d}t \int_{\Sigma} \Bigg[
    e_{0}{}^{a} \, \epsilon^{ij} \Big(
      \partial_{i}\omega_{j}^{a}
      - \partial_{j}\omega_{i}^{a}
      + \epsilon^{abc}\omega_{ib}\omega_{jc}
      + \Lambda \, \epsilon^{abc} e_{ib} e_{jc}
    \Big) \\
    & \qquad\qquad
    + \, \omega_{0}{}^{a} \, \epsilon^{ij} \Big(
      \partial_{i}e_{j}^{a}
      - \partial_{j}e_{i}^{a}
      + \epsilon^{abc}\big(
        \omega_{ib} e_{jc}
        + e_{ib} \omega_{jc}
      \big)
    \Big)
  \Bigg] \, .
\end{align}
\rmk{
  Here $ i,j = 1,2 $ which are the spacial index. Thus we can view $
  e_i{}^a $ and $ \omega_i{}^a $ as some 1-forms defined on the
  spacial slice $ \Sigma_g $.
}
By looking at the Lagrangian, we can see many properties:
\begin{itemize}
  \item \textbf{Constraints} $ e_0{}^a $ and $ \omega_0{}^a $ are not
    dynamical variables but lagrangian mulipliers that gives us the
    constraints.

    We can get the constriant equation by varying the
    action with respect to $ e_0{}^a $ and $ \omega_0{}^a $:
    \begin{align}
      \epsilon^{ij} \Big(
        \partial_{i}\omega_{j}^{a}
        - \partial_{j}\omega_{i}^{a}
        + \epsilon^{abc}\omega_{ib}\omega_{jc}
        + \Lambda \, \epsilon^{abc} e_{ib} e_{jc}
      \Big) & = 0 \\
      \epsilon^{ij} \Big(
        \partial_{i}e_{j}^{a}
        - \partial_{j}e_{i}^{a}
        + \epsilon^{abc}\big(
          \omega_{ib} e_{jc}
          + e_{ib} \omega_{jc}
        \big)
      \Big) & = 0
    \end{align}
    We can see that compare to \cref{eq:curvaturetensorcos} that the
    constraint equation are the spacial part of the curvature tensor
    $ F_{ij} = 0, i,j = 1,2 $. the constiants equation shows us that
    $ A_i $ on the spacial slice $ \Sigma_g $ is a flat connection.
  \item we can see that the $ \omega_i{}^a $ and $ e_i{}^a $ on the
    spacial slice $ \Sigma $ are the canonical variables and being
    the canonical conjugate of each other. Thus we can define a
    classical phase space and Poisson Bracket on it.
\end{itemize}
\thm{
  Pre-Constraint Phase Space and Poisson Bracket

  The Pre-Constraint Phase Space is space of configuration of $
  ( e_i{}^a , \omega_i{}^a) $ on the spacial slice $ \Sigma_g
  $.(resembling $ (p,q) $ configuration on a equal time slice in the
  standard Hamiltonian formalism)

  The Poisson Bracket is defined as:
  \begin{align}
    & \{ e_i{}^a(x), \omega_j{}^b(y) \} = \frac{1}{2} \epsilon_{ij} \eta^{ab}
    \delta^{(2)}(x-y)\\
    &\{ e_i{}^a(x), e_j{}^b(y) \} = 0 \\
    &\{ \omega_i{}^a(x), \omega_j{}^b(y) \} = 0
  \end{align}
  This definition is consistent with the Lagrangian.
}

\subsubsection{Constraint Quantization and Phase Space}

In Dirac quantization, we use the pre-constraint phase space to
quantize the theory, and view constriant as operator. The Hilbert
space is the space of states that are annihilated by the constraint
operator.

Here we would prefer to use another procedure, called the "constraint
first quantization". The idea is that find a "True Phase Space" with
the following steps and quantize the phase space directly:
\begin{enumerate}
  \item \textbf{Find Solution Space of Constraint: } given may
    constraints $ H_I, I = 1,2,3 \dots $ find the solution space of
    the constraint equations $ H_I = 0 $. This is a subspace of the
    pre-constraint phase space.
  \item \textbf{Find Equivalence Class: } the constraint $ H_I $
    may generate a transformation on the canonical vairable $ (p,q) $
    using the Poisson Bracket:
    \begin{align}
      \delta_C q = \displaystyle\sum_{I} \epsilon_I \{ H_I, q \}
      \quad \delta_C p = \displaystyle\sum_{I} \epsilon_I \{ H_I, p \}
    \end{align}
    We need to find the equivalence class of the solution space under
    this transformation, which is called the "True Phase Space". Or
    equivalently, the space of transformation invariant function on
    the solution space.
  \item \textbf{Poisson Bracket on True Phase Space:} Poisson Bracket
    is unchanged, long as we only consider the transformation
    invariant function on the solution space, which is the "True Phase Space".
\end{enumerate}
For our case of 2+1 dimensional gravity. The solution space of the
constriant is the space of flat $ G \in \{ ISO(2,1), SO(3,1),
SO(2,2)\} $ connection on the spacial slice $ \Sigma_g $.
\rmk{
  we need to be careful to distinguish the phase space with the EoM
  solution space. The EoM solution space is the space of flat G
  connection on the whole spacetime. The phase space is the space of flat G
  connection on the spacial slice.
}
And amazingly we can find that the transformation generated by the
constraint is just the $ G $ gauge transformation on the spacial
slice \cref{eq:gaugetransformcos}.
\rmk{
  In fact in gauge theory, the constraint always generated the gauge
  transformation.
}
Thus we have the result:
\thm{\label{thm:True Phase Space of 2+1 Dimensional Gravity}
  True Phase Space of 2+1 Dimensional Gravity

  The True Phase Space of 2+1 dimensional gravity is the moduli space of
  flat $ G \in \{ ISO(2,1), SO(3,1), SO(2,2)\} $ connection on the
  spacial slice $ \Sigma_g $. Mathematically, this is equivalent to
  the following space:
  \begin{align}
    \mathcal{M} = \mathrm{Hom}(\pi_1(\Sigma_g), G) / \sim
  \end{align}
  where $ \sim $ is the equivalence relation given by the conjugation of
  $ G $ of group homomorphism $ \mathrm{Hom}(\pi_1(\Sigma_g), G) $.
}

\subsubsection{Compare with Geometric Analysis of Phase
Space}\label{sec:comparephase}

We now construct a "True Phase Space" using the constraint first
approach of the Chern-Simons formulation of 2+1 dimensional gravity
\cref{thm:True Phase Space of 2+1 Dimensional Gravity}. We can
compare it with the geometric analysis of the phase space we've done
initially \cref{thm:Geophasespace}.
We can see that the result is \textbf{Exactly the Same!!!!}, apart
form using different approaches.

However, we shall remark on something we have negelected which make
2+1 dimensional gravity not exactly equivalent to Chern-Simons Theory:
\begin{itemize}
  \item Diffeomorphism is not exactly the gauge transformations.
    Gauge transformation only covers "small diffeomorphism", while
    "large diffeomorphism" is not included in the gauge
    transformation.
  \item In 2+1 Dimensional Gravity we assume that $ e_i{}^a $ is
    invertible and want the metric to behave well. While the
    Chern-Simons theory doesn't care about this.
\end{itemize}

\subsection{Real Polarization Quantization}

Now we focus on quantizing the "True Phase Space". There are many
approaches to quantize a phase space.
\begin{itemize}
  \item If the phase space $ \mathcal{M} $ is a cotangent bundle of
    some configuration space $N$ then we can use a approach called
    \textbf{Real Polarization}
\end{itemize}
This is the true for Minkowski and dS case, but not for AdS case.
This section we'll focus on the Minkowski case.

\subsubsection{True Phase Space as Cotangent Bundle}

A navie idea when looking at the pre-constraint phase space $
(e_i{}^a, \omega_i{}^a) $ is that we can view $ e_i{}^a $ as the
momentum and $ \omega_i{}^a $ as the coordinate, thus the phase space
is a cotangent bundle of the space of $ \omega_i{}^a $ configuration.

But this is note generally true. The problem is that the constraint
generate a transformation that mixed $ e_i{}^\mu $ and $ \omega_i{}^a
$. Thus we can't simply view $ e_i{}^a $ as the momentum and $
\omega_i{}^a $ as the coordinate.

However, if we focus on the case that $ \Lambda = 0 $, we may see
from \cref{eq:gaugetransformcos} that the gauge transformation is
purely a transformation on $ \omega_i{}^a $ and doesn't mix with $
e_i{}^a $. Thus:
\begin{itemize}
  \item for $ \Lambda = 0 $ case admits a real polarization
    quantization, where we can view $ e_i{}^a
    $ as the momentum and $ \omega_i{}^a $ as the coordinate, thus the
    phase space is a cotangent bundle of the space of $ \omega_i{}^a $
    configuration moduli the gauge transformation.
\end{itemize}
We notice that $ \omega_i{}^a $ itself is infact a $ SO(2,1) $ gauge
field on the spacial slice $ \Sigma_g $. Thus we have the conclusion that:
\thm{
  True Phase Space as Cotangent Bundle $ \Lambda = 0 $

  The true phase space of 2+1 dimensional gravity with $ \Lambda = 0
  $ is a cotangent bundle with $ \mathcal{M} = T^* N $, where $ N $ is
  the space of flat $ SO(2,1) \sim SL(2,\mathbb{R}) $ connection on
  the spacial slice $ \Sigma_g $.
  \begin{align}
    \mathcal{M} = T^* N  \quad N = Hom(\pi_1(\Sigma_g),
    SL(2,\mathbb{R})) / \sim
  \end{align}
}
\rmk{
  This is not a rigorous proof, but it have intuitively give us a
  picture of the phase space.
}

\subsubsection{Hilbert Space}

Then as a analogue to the standard quantum mechanics. The Hilbert
space can be view as a function space (all possible wave function) on
the configuration space $ N $:
\begin{align}
  \psi(x) \in \mathcal{H} \quad x \in N
\end{align}
and it shall be square integrable:
\begin{align}
  \int_{N} |\phi(x)|^2 d \mu< \infty
\end{align}
where $ d \mu $ is a suitable measure on the configuration space $ N
$. However, in many cases, there isn't a natural measure on the
configuration space, thus we shall define the element of Hilbert
space not as measure but some "half density" on the configuration
space, which is a sort of a square root of a volume element.

To explicitly construct the Hilbert Space, we can label the element
of the configuration space $ N $ by some coordinates $ U^i,V^i \in
SL(2,\mathbb{R}) $.

This is because a homomorphism $ \mathrm{Hom}(\pi_1(\Sigma_g),
SL(2,\mathbb{R})) $ is given by a set of image of the generators of $
\pi_1(\Sigma_g) $ which are  $ \phi(a^i,b^i) = U^i,V^i \in
SL(2,\mathbb{R}) $ that satisfy the relation:
\begin{align}
  \prod_{i=1}^g U^i V^i (U^i)^{-1} (V^i)^{-1} = 1
\end{align}
Thus we can label the element of the configuration space $ N $ by $
U^i,V^i \in SL(2,\mathbb{R}) $ that satisfy the above relation. Then
the Hilbert space can be view as a function space on $ U^i,V^i $:
\begin{align}
  \phi(U^i,V^i) \in \mathcal{H}
\end{align}
That satisfy:
\begin{enumerate}
  \item Being invariant under the gauge transformation, which is the
    conjugation of
    $ SL(2,\mathbb{R}) $:
    \begin{align}
      \psi(U^i,V^i) = \psi(g U^i g^{-1}, g V^i g^{-1})
    \end{align}
  \item Being square integrable:
    \begin{align}
      |\psi|^2 = \int \psi(U^i,V^i) * \psi(U^i, V^i)  < \infty
    \end{align}
\end{enumerate}

\subsubsection{Topological of Flat $ SL(2,\mathbb{R}) $
Connection}\label{sec:Topologicalsl2r}

If we investigate the configuration space $ N $ carefully. We may
find that this is not a connected space, it can divided into several
connected components. Let's investigate the classification and
properties of these connected components.

To see his classification, we can first notice that a $
SL(2,\mathbb{R}) $ connection on $ \Sigma_g $ can induce a real plane
bundle on $ \Sigma_g $. The idea is that the element of $
SL(2,\mathbb{R}) $ can be view as a linear transformation on a real plane.
\YL{[It may take more mathematical technique to see this, we'll not do it here]}
For flat $ SL(2,\mathbb{R}) $ connection gives out a flat plane
bundle, which is classified by the Euler class $ e \in [-(2g-2),
2g-2] $ of the plane bundle. The configuration space $ N $ can be
divided into several connected components $ N_e $ labeled by the
Euler class $ e $ of the plane bundle.

Now the problem is that which connected component do we want? and why??
\begin{itemize}
  \item The answer is that we want the component with $ e = 2g-2 $ !!
\end{itemize}
The reason is due to an important mathematical fact:
\thm{
  The space of homomorphism $ \phi: \pi_1(\Sigma_g) \to
  SL(2,\mathbb{R}) $ where the image $ \phi(\pi_1(\Sigma)) $ is a
  discrete subgroup of $ SL(2,\mathbb{R}) $ is exactly the connected
  component $ N_e $ with $ e = 2g - 2 $.
  \begin{align}
    \text{Discrete Enbeding of } Hom(\pi_1(\Sigma_g),
    SL(2,\mathbb{R}))/sim = N_{e=2-2g}
  \end{align}
}
And we want the discrete embedding! An geometrically argument is that
it may lead to a well-behaved quotient space $ \mathbb{H}^2 /
\phi(\pi_1(\Sigma_g)) $ which is a hyperbolic surface.

\YL{[in fact I don't understand why we want this discrete embedding.
Perhaps it is just a guess of Witten]}

Morover, we know that a discrete subgroup of $ \phi(\pi_1(\Sigma_g) = \Gamma
  \subset SL(2,\mathbb{R}) $
  which is a discrete subset of isometry group of $ \mathbb{H}^2 $,
  and each one of them may lead to a manifold with the topology of $
  \Sigma_g $ by doing a quotient of $ \mathbb{H}^2 $ by the discrete
  subgroup $ \Gamma $. The space of all these manifolds is exactly the
  Teichmuller space $ \mathcal{T}(\Sigma_g) $ of $ \Sigma_g $. Thus
  we have the result that this component $ N_{e=2-2g} $ is the
  Teichmuller space $ \mathcal{T}(\Sigma_g) $.

  \subsection{Kahler Polarization Quantization}
  For AdS case, the real polarization quantization can't work, thus we
  can follow another procedure that may work for both AdS, Minkowski
  and dS case, which is the Kahler polarization quantization.
  The idea is that:
  \begin{itemize}
    \item First we find a Kahler structure on the phase space.
    \item Then we find a holomorphic sector of a line bundle over the
      phase space and this may give us a Hilbert space.
  \end{itemize}
  This procedure is not discussed in Witten's paper, but it is widely
  used in the modern literature.
